\documentclass[11pt]{article}
\usepackage{amsmath,amssymb,amsthm}
\newtheorem{proposition}{Proposition}
\title{A 2-adic point on the five-row comparison variety}
\author{Exact computational certificate; independently reviewed}
\date{30 September 2026}
\begin{document}
\maketitle
\begin{proposition}
The five-row, four-column centered comparison moment equations have
a solution over $\mathbb Q_2$. Thus 2-adic local insolubility cannot
exclude rational five-row models.
\end{proposition}
\begin{proof}
Write $\mathbb E$ for averaging five coordinates. Set
\[
 r=3/4,\quad s=1/4,\quad p=C/16,\quad q=5/6-p,
 \quad R=p/r,\quad S=q/s,
\]
and put
\[
 D=r^2q-s^2p,\quad
 z_1=\frac{5r(q-s^2)}{6D},\qquad
 z_2=\frac{5s(r^2-p)}{6D}.
\]
Let
\[
 z=(-z_1,-z_2,0,z_2,z_1),\qquad
 u=(q,-p,2(p-q),-p,q).
\]
We seek the four columns
\[
 a=z+\lambda u,\quad b=z-\lambda u,\quad
 c=(-r,-s,0,s,r),\quad d=(-R,-S,0,S,R).
\]
The means vanish. Direct substitution gives
$\mathbb Ecd=\mathbb Ezc=\mathbb Ezd=1/3$ and
$\mathbb Eucd=0$. Oddness under row reversal then verifies all
four triple moments and the other mixed pair moments. The pair $ab$
has the correct moment when
\[
 \lambda^2=L(C):=
 \frac{64(C^2-24C+148)}{3(C-12)^2(3C^2-40C+160)}.
\]
Under this condition the remaining fourth moment
$\mathbb Eabcd=1/5$ is equivalent, whenever the displayed
denominators are nonzero, to
\[
 f(C):=36C^5-1647C^4+29612C^3-264376C^2+1193920C-2229120=0.
\]
All these identities can also be checked by clearing their rational
polynomial denominators.

Take $C_0=2179188$. Exact integer evaluations give
\[
 v_2(f(C_0))=28,\qquad v_2(f'(C_0))=6.
\]
Hensel's lemma therefore gives a root $C_*\in\mathbb Z_2$ with
$C_*-C_0\in2^{22}\mathbb Z_2$. For the displayed integer numerator
$P(C)$ and denominator $Q(C)$ of $L(C)$,
\[
 v_2(P(C_0))=8,\qquad v_2(Q(C_0))=10,
 \qquad 4L(C_0)\equiv1\pmod8.
\]
Integer polynomials change by a multiple of $2^{22}$ throughout
this ball. It follows that $v_2(L(C_*))=-2$ and
$4L(C_*)\equiv1\pmod8$. The usual square criterion in $\mathbb Q_2$
now supplies a square root $\lambda\in\mathbb Q_2$.
The same valuation calculation keeps $Q(C_*)$ nonzero. Since
\[
 D=-5(C-12)/128,\qquad
 \mathbb Eu^2=5/6-4pq=(3C^2-40C+160)/192,
\]
all denominators needed in the construction are nonzero. The four
columns therefore satisfy all fifteen required moment equations.
\end{proof}
This is a local-solubility result only. It proves neither a rational
model nor a monotone real model, and makes no claim of realizability
by actual dice. An elementary separate argument shows that the
centered coordinates of any such model cannot all lie in $\mathbb Z_2$;
the point here indeed has negative valuations.
\end{document}
